\documentclass[12pt,a4paper]{article}
\usepackage{amsmath,amssymb,amsthm}
\usepackage{booktabs}
\usepackage{enumitem}
\usepackage{hyperref}
\usepackage{geometry}
\geometry{margin=2.5cm}

\newtheorem{theorem}{Theorem}[section]
\newtheorem{lemma}[theorem]{Lemma}
\newtheorem{corollary}[theorem]{Corollary}
\newtheorem{proposition}[theorem]{Proposition}
\theoremstyle{definition}
\newtheorem{definition}[theorem]{Definition}
\theoremstyle{remark}
\newtheorem{remark}[theorem]{Remark}

\title{The Speed of Light from Recognition Science:\\
One Voxel per Tick as Structural Speed Limit}

\author{Jonathan Washburn\\
Recognition Science Research Institute\\
Austin, Texas\\
\texttt{washburn.jonathan@gmail.com}}

\date{March 2026}

\begin{document}
\maketitle

\begin{abstract}
We derive the speed of light $c$ from the discrete structure of the
recognition ledger.  The ledger consists of voxels of size $\ell_0$
updated at tick intervals $\tau_0$.  The maximum propagation speed is
one voxel per tick: $c = \ell_0/\tau_0$.  In RS-native units
($\ell_0 = \tau_0 = 1$), $c = 1$ exactly.  The light-cone identity
$c\,\tau_0 = \ell_0$ is a structural consequence, not a coincidence.
The constancy of $c$ (its independence from source motion) follows from
the fact that the lattice tick rate is universal.  The SI value
$c = 299\,792\,458\;\mathrm{m/s}$ is recovered through the calibration
seam (identification of $\ell_0 = \ell_P$, $\tau_0 = t_P$).
\end{abstract}

\section{Introduction}
\label{sec:intro}

The speed of light $c$ is the universal speed limit in physics.
Einstein's 1905 postulate~\cite{Einstein1905}---that $c$ is the
same in all inertial frames---is confirmed to extraordinary
precision, most dramatically by the Michelson--Morley null
result~\cite{MichelsonMorley1887} and the coincident detection
of gravitational and electromagnetic signals from
GW170817~\cite{GW170817}.  Yet special relativity offers no
explanation for \emph{why} there is a universal speed limit or
why it takes the value it does.

Recognition Science~\cite{RS_Axioms} derives $c$ from the discrete
ledger.  The ledger is a $D = 3$ cubic lattice of voxels with
spatial quantum $\ell_0$ and tick duration $\tau_0$.  Information
propagates at most one voxel per tick; massless signals saturate
this bound.  The speed of light is therefore
\begin{equation}
  c = \frac{\ell_0}{\tau_0} = 1\quad\text{(RS-native units)}.
\end{equation}
Its constancy, its role as a speed limit, and the Lorentz group
all follow from the lattice structure.

\section{Recognition Science Framework}
\label{sec:framework}

Recognition Science derives all physics from a single primitive,
the \emph{recognition cost functional} $J(x)$, uniquely determined
by three axioms.

\begin{definition}[RS Cost Functional]
For $x > 0$:
$J(x) = \tfrac{1}{2}(x + x^{-1}) - 1$.
\end{definition}

\noindent\textbf{A1 (Normalization):} $J(1) = 0$.\quad
\textbf{A2 (RCL):} $J(xy) + J(x/y) = 2J(x)J(y) + 2J(x) + 2J(y)$.\quad
\textbf{A3 (Calibration):} $J''_{\log}(0) = 1$.

\begin{theorem}[Cost Uniqueness~\cite{Aczel1966}]
\label{thm:unique}
The continuous solution to \textup{(A1)--(A3)} is $J(x) =
\tfrac{1}{2}(x + x^{-1}) - 1$ (unique by the Acz\'el classification of
d'Alembert's functional equation).
\end{theorem}

\noindent The golden ratio $\varphi = (1+\sqrt{5})/2$ is forced by
self-similarity.  Spatial dimension $D = 3$ is forced by the linking
condition $\operatorname{lcm}(2^D,45) = 360$: the 8-tick
epoch ($2^3 = 8$) and the angular quantum ($45^\circ$) tile the
full circle ($360^\circ$) only when $D = 3$.
The $J$-cost depends only on ratios, giving the ledger its
translation-invariant and Lorentz-compatible structure.

\section{The Ledger Structure}

\begin{definition}[Fundamental Scales]
The recognition ledger has:
\begin{itemize}
  \item Voxel size $\ell_0 > 0$ (the minimum resolvable length).
  \item Tick duration $\tau_0 > 0$ (the minimum resolvable time).
\end{itemize}
Both are forced by the 8-tick epoch: $\tau_0$ is one tick of the
$2^D = 8$ cycle ($D = 3$); $\ell_0$ is the corresponding spatial
quantum.
\end{definition}

\section{Derivation}

\begin{proposition}[Speed of Light --- Structural Derivation]
\begin{equation}
  \boxed{c = \frac{\ell_0}{\tau_0}.}
\end{equation}
In RS-native units ($\ell_0 = \tau_0 = 1$): $c = 1$.
\end{proposition}

\begin{proof}
The ledger is a discrete lattice: voxels at integer coordinates
$(i,j,k) \in \mathbb{Z}^3$ updated at integer ticks $t \in
\mathbb{Z}$.  A recognition event at voxel $\mathbf{x} = (i,j,k)$ at
tick $t$ can influence only those voxels that share an edge, face, or
vertex with $\mathbf{x}$ at tick $t+1$---the nearest neighbors in the
$3+1$ causal graph.  The maximum spatial displacement in one tick is
one voxel: $\|\Delta\mathbf{x}\| \leq \ell_0$ when $\Delta t =
\tau_0$.  Therefore
\begin{equation}
  v_{\max} = \frac{\|\Delta\mathbf{x}\|_{\max}}{\Delta t}
  = \frac{\ell_0}{\tau_0}.
\end{equation}
Massless signals (electromagnetic and gravitational) propagate by
nearest-neighbor ledger updates; they saturate this bound.  Hence
$c = \ell_0/\tau_0$.  The bound is structural: no physical process
can propagate faster without violating the discrete causal structure.
\end{proof}

\begin{corollary}[Light-Cone Identity]
$c\,\tau_0 = \ell_0$, and $c > 0$.
\end{corollary}

\section{Why $c$ Is Constant}
\label{sec:constant}

\begin{proposition}[Frame Independence --- Structural Argument]
\label{thm:constant}
The speed $c = \ell_0/\tau_0$ is independent of the motion of
the source or the observer.
\end{proposition}

\begin{proof}
Three properties of the ledger force this:
\begin{enumerate}[label=(\roman*),nosep]
  \item \emph{Universal tick rate.}  There is one ledger, not
  many.  The tick rate $\tau_0$ is a global property of the
  ledger, not a local property of any source or observer.
  \item \emph{Homogeneity.}  The voxel size $\ell_0$ is
  uniform---the ledger has no distinguished location because the
  $J$-cost depends only on ratios, not absolute positions.
  \item \emph{Cost invariance.}  The $J$-cost functional
  $J(x) = \frac{1}{2}(x+x^{-1})-1$ depends only on ratios,
  so the dynamics are invariant under lattice translations.
\end{enumerate}
Since both $\ell_0$ and $\tau_0$ are universal, their ratio is
the same for all observers:
$c = \ell_0/\tau_0$ regardless of relative motion.
\end{proof}

\begin{remark}
In special relativity, the constancy of $c$ is an axiom.  In RS,
it is a theorem: a consequence of the ledger having one global
tick rate and uniform voxel size.  The Michelson--Morley null
result~\cite{MichelsonMorley1887} is explained structurally:
there is no ``ether'' to produce directional variation because
the lattice is isotropic (Proposition~\ref{prop:isotropy}).
\end{remark}

\section{Lorentz Invariance}
\label{sec:lorentz}

\begin{proposition}[Lattice Isotropy]
\label{prop:isotropy}
The $D = 3$ cubic lattice has discrete rotation symmetry
(the octahedral group $O_h$, 48 elements).  In the continuum
limit $\ell_0 \to 0$, this discrete symmetry is promoted to
full $SO(3)$ rotational invariance.
\end{proposition}

\begin{proposition}[Minkowski Metric from the Lattice --- Continuum Limit]
\label{thm:minkowski}
In the continuum limit $\ell_0, \tau_0 \to 0$ with $c = \ell_0/\tau_0$
fixed, the discrete causal structure of the ledger becomes the
invariant interval
\begin{equation}
  ds^2 = c^2\,dt^2 - dx^2 - dy^2 - dz^2.
\end{equation}
\end{proposition}

\begin{proof}
On the lattice, causal separation is classified by the number
of ticks and voxel steps between two events.
Nearest-neighbor propagation at one voxel per tick gives
$c\,\Delta t = \|\Delta\mathbf{x}\|$, i.e.,
$ds^2 = 0$ (null separation).  Events with $ds^2 > 0$ are
timelike (connected by sub-light paths); events with
$ds^2 < 0$ are spacelike (no causal connection).  The discrete
light cone---all voxels reachable in exactly $n$ ticks---defines
the null surface.  In the continuum limit
$\ell_0, \tau_0 \to 0$ with $c = \ell_0/\tau_0$ fixed, the
discrete rotational symmetry ($O_h$) is promoted to $SO(3)$,
and the lattice causal structure becomes the Minkowski metric
$\eta_{\mu\nu} = \mathrm{diag}(+1,-1,-1,-1)$.
\end{proof}

\begin{remark}[Non-triviality of the continuum limit]
The promotion from discrete $O_h$ (octahedral, 48 elements) to
continuous $SO(3,1)$ (Lorentz group) in the continuum limit is
non-trivial.  The cubic lattice has 48 rotational symmetries, not
the full $SO(3)$.  At the level of the $J$-cost functional (which
depends only on ratios, not absolute directions), the lattice lacks
a preferred frame, and the Lorentz group emerges as the symmetry of
the causal cone in the limit $\ell_0 \to 0$.  However, a complete
proof that discrete $O_h$ is promoted to continuous $SO(3,1)$ in this
limit requires a careful treatment of how lattice symmetry is enhanced
at long wavelengths.  This is an identified open problem in the RS
formalism.
\end{remark}

\begin{remark}
A discrete lattice naively breaks continuous Lorentz invariance.
RS addresses this in two ways:
\begin{enumerate}[label=(\roman*),nosep]
  \item The $J$-cost functional $J(x) = \frac{1}{2}(x+x^{-1})-1$
  depends only on ratios, not absolute positions---the lattice
  has no preferred ``rest frame'' at the cost level.
  \item In the continuum limit, all lattice artifacts vanish at
  scales $\lambda \gg \ell_0$; deviations from Lorentz symmetry
  are suppressed by $(\ell_0/\lambda)^2$.
\end{enumerate}
This is analogous to how lattice QCD preserves gauge invariance
despite working on a discrete grid.
\end{remark}

\section{Why $c$ Cannot Be Exceeded}

\begin{proposition}[Speed Limit from $J$-Cost --- Structural Argument]
No signal can propagate faster than $c = \ell_0/\tau_0$.
\end{proposition}

\begin{proof}
Suppose a signal propagated at speed $v > c$.  In one tick $\tau_0$, it
would traverse a distance $v\,\tau_0 > \ell_0$, skipping over at least
one voxel.  The ledger would then record an influence at a voxel
$\mathbf{x}'$ that is not a nearest neighbor of the source voxel
$\mathbf{x}$ at the previous tick.  Such a ``jump'' creates an
unbalanced ledger entry: a debit (or credit) at $\mathbf{x}'$ without
a corresponding credit (or debit) at an intermediate voxel along the
path.  By the double-entry axiom of the recognition ledger, unbalanced
entries have infinite $J$-cost.  Cost-minimizing dynamics therefore
exclude any propagation with $v > c$.  Hence $v \leq c$ for all
physical signals.
\end{proof}

\begin{corollary}
Massive particles have $v < c$ because they require multiple ticks per
voxel step (they do not saturate the light cone).  Only massless
signals achieve $v = c$.
\end{corollary}

\section{Gravitational Waves at $c$}

Gravitational waves propagate at exactly $c = \ell_0/\tau_0$.  Both
EM and GW signals are massless excitations of the same discrete
ledger, so both saturate the structural speed bound.

\begin{proposition}[GW Propagation Speed --- Structural Argument]
$c_{\mathrm{gw}} = c = \ell_0/\tau_0$.
\end{proposition}

\begin{proof}[Structural argument]
The RS ledger has one causal cone: one voxel per tick.  Both
electromagnetic and gravitational perturbations propagate on the same
lattice with the same voxel size $\ell_0$ and tick duration $\tau_0$.
The key RS claim is that the same causal structure governs all massless
excitations, so $c_{\mathrm{gw}} = \ell_0/\tau_0 = c$.

\emph{Note on circularity}: The wave equation $\Box h_{\mu\nu} = 0$
with $\Box = \partial_t^2 - c^2\nabla^2$ has $c$ built in by definition,
making it appear circular to conclude $c_{\mathrm{gw}} = c$ from this
equation.  The non-circular content of the RS argument is that the
d'Alembertian is defined by the \emph{same} lattice parameters
($\ell_0$, $\tau_0$) for gravitons as for photons---i.e., gravity
and electromagnetism propagate on the same substrate, not different
substrates with different speeds.
\end{proof}

\begin{remark}
The argument just given is the RS structural proof that $c_{\mathrm{gw}} = c$
is a \emph{non-trivial identity}, not a definitional tautology: both
quantities are independently defined (the ledger causal speed and the
GW propagation speed), and their equality follows from the single-substrate
lattice structure.  See the companion paper on GW propagation for a
more detailed treatment.
\end{remark}

\begin{remark}
GW170817 (binary neutron star merger with coincident gamma-ray burst
GRB~170817A) confirmed $|c_{\mathrm{gw}}/c - 1| < 10^{-15}$~\cite{GW170817}.
RS predicts the ratio is exactly 1---a structural identity, not an
approximation.
\end{remark}

\section{Connection to Other Constants}
\label{sec:constants}

The speed of light sits at the head of the RS derivation chain:
\begin{equation}
  \mathrm{RCL} \;\to\; J \;\to\; \varphi \;\to\;
  \tau_0\;\text{(8-tick)} \;\to\;
  \boxed{c = \frac{\ell_0}{\tau_0} = 1} \;\to\;
  \hbar = \varphi^{-5} \;\to\; G = \varphi^5.
\end{equation}
Once $c$ is established, the remaining fundamental constants follow:

\begin{proposition}[Planck Length Identity]
\label{prop:planck}
In RS-native units ($c = 1$, $\hbar = \varphi^{-5}$,
$G = \varphi^5$):
\begin{equation}
  \ell_P = \sqrt{\frac{\hbar G}{c^3}}
  = \sqrt{\frac{\varphi^{-5} \cdot \varphi^5}{1}}
  = 1.
\end{equation}
\end{proposition}

\begin{remark}
The Planck length equals the voxel size, and the Planck time
equals the tick duration.  There is no hierarchy between the lattice
scale and the Planck scale---they are the same scale.  This is a
structural identity, not a numerical coincidence.
\end{remark}

\begin{remark}[Convention dependence of $G$]
\label{rem:G-convention}
The identity $\ell_P = 1$ depends on the normalization of $G$.
The formula displayed above uses $G = \varphi^5$ in RS-native
units ($c = \hbar = 1$), which gives $\ell_P = 1$ exactly.
An alternative RS derivation yields $G = \varphi^5/\pi$ (with
the factor of $\pi$ absorbed into the recognition wavelength
convention); under that convention $\ell_P = 1/\sqrt{\pi}$
and the voxel--Planck identification acquires a geometric
prefactor.  The precise $G$ formula and its interplay with the
Planck scale are treated in the companion paper on the
gravitational constant.
\end{remark}

\section{SI Calibration}
\label{sec:si}

The calibration seam maps RS-native units to SI by identifying
$\ell_0$ with the Planck length and $\tau_0$ with the Planck time.
Since $\ell_P = \ell_0 = 1$ in RS-native units:
\begin{align}
  \ell_0 &= \ell_P \approx 1.616 \times 10^{-35}\;\mathrm{m}, \\
  \tau_0 &= t_P \approx 5.39 \times 10^{-44}\;\mathrm{s}.
\end{align}
The SI speed of light is then:
\begin{equation}
  c_{\mathrm{SI}} = \frac{\ell_P}{t_P}
  = \frac{1.616 \times 10^{-35}}{5.39 \times 10^{-44}}
  \approx 2.998 \times 10^8\;\mathrm{m/s}.
\end{equation}
This reproduces the exact SI definition
$c = 299\,792\,458\;\mathrm{m/s}$.  RS does not predict this
number---it is a unit convention.  What RS predicts is that $c = 1$
in natural units, i.e., the speed of light equals one voxel per tick.

\section{No-Dispersion Theorem}
\label{sec:dispersion}

Some discrete-spacetime models predict energy-dependent photon
velocities.  RS does not.

\begin{proposition}[No Dispersion --- Structural Prediction]
\label{thm:nodispersion}
All massless signals propagate at exactly $c = \ell_0/\tau_0$,
independent of wavelength, polarization, or energy.
\end{proposition}

\begin{proof}
The bound $v \leq \ell_0/\tau_0$ is purely structural: it
depends only on the nearest-neighbor topology of the lattice,
not on the energy or frequency of the signal.  A massless signal
saturates the bound regardless of its wavelength $\lambda$
(provided $\lambda \geq \ell_0$, which is always true below
the Planck energy).  No energy-dependent correction to $c$
exists because no RS mechanism introduces one: the $J$-cost
of a nearest-neighbor update is the same for all massless modes.
\end{proof}

\begin{remark}
This prediction is testable.  Fermi-LAT observations of
GRB~090510 constrain Lorentz-violating dispersion to
$|v(E)/c - 1| < 10^{-20}$ at photon energies
$E \sim 30\;\mathrm{GeV}$~\cite{FermiLAT2009}.  RS predicts
exact equality $v(E) = c$ at all energies---a stronger claim
than current bounds can test.
\end{remark}

\section{Predictions}
\label{sec:predictions}

\begin{enumerate}
  \item \textbf{Exact $c_{\mathrm{gw}}/c = 1$.}  Gravitational and
  electromagnetic waves propagate at identical speeds.  Any future
  measurement showing $|c_{\mathrm{gw}}/c - 1| > 10^{-20}$ would
  falsify RS.

  \item \textbf{No superluminal signaling.}  No physical process
  (including quantum entanglement) can transmit information faster than
  $c$.  The ledger update propagates at one voxel per tick; no
  mechanism can skip voxels without infinite $J$-cost.

  \item \textbf{Zero photon dispersion.}  All massless fields
  propagate at exactly $c$, independent of energy.  RS predicts
  no Lorentz-violating dispersion at any scale, including the
  Planck scale (Theorem~\ref{thm:nodispersion}).

  \item \textbf{Planck-scale discrete light cone.}  At the Planck
  scale, the light cone is a discrete lattice surface, not a smooth
  null cone.  Continuum relativity emerges from coarse-graining at
  $\lambda \gg \ell_0 = \ell_P$.

  \item \textbf{No massive photon.}  The photon mass is exactly
  zero: photons saturate the causal bound.  Any nonzero mass
  would yield $v < c$, contradicting the structural identity
  $v_{\gamma} = \ell_0/\tau_0$.
\end{enumerate}

\section{Comparison}

\begin{center}
\renewcommand{\arraystretch}{1.3}
\begin{tabular}{p{3.2cm}p{5.5cm}p{5.5cm}}
\toprule
\textbf{Feature} & \textbf{Special Relativity} & \textbf{Recognition Science} \\
\midrule
Origin of $c$ & Postulate (invariant speed) & Derivation ($\ell_0/\tau_0$) \\
Why $c$ is constant & Axiom & Lattice universality \\
Speed limit & Postulate (no $v > c$) & $J$-cost (unbalanced entries forbidden) \\
Minkowski metric & Postulate & Emerges from discrete light cone \\
$c_{\mathrm{gw}} = c$ & Assumed in GR & Structural identity \\
Lorentz group & Symmetry of spacetime & Symmetry of causal lattice \\
\bottomrule
\end{tabular}
\end{center}

\section{Conclusion}

The speed of light is derived from the discrete ledger:
$c = \ell_0/\tau_0 = 1$ voxel per tick.  Its constancy follows
from lattice universality; it cannot be exceeded because
superluminal propagation incurs infinite $J$-cost.  The Minkowski
metric emerges from the discrete causal structure in the continuum
limit; the Lorentz group is the symmetry of the causal lattice.
No dispersion or energy-dependent deviation from $c$ exists in RS.

The connection to other constants is structural: $\ell_0 = \ell_P$,
$\tau_0 = t_P$, and the Planck length identity $\ell_P = 1$ in
RS-native units follows from $G \cdot \hbar = 1$.  The speed of
light is not a free parameter but the first link in the derivation
chain from $\varphi$ to all physical constants.

\begin{thebibliography}{99}

\bibitem{Aczel1966}
J.~Acz\'el, \emph{Lectures on Functional Equations and Their Applications},
Academic Press, New York (1966).

\bibitem{Einstein1905}
A.~Einstein, ``On the Electrodynamics of Moving Bodies,'' Ann.\
Phys.\ \textbf{17}, 891 (1905).

\bibitem{MichelsonMorley1887}
A.~A.~Michelson and E.~W.~Morley, ``On the Relative Motion of the
Earth and the Luminiferous Ether,'' Am.\ J.\ Sci.\
\textbf{34}, 333 (1887).

\bibitem{GW170817}
LIGO/Virgo Collaboration, ``Gravitational Waves and Gamma-Rays from
a Binary Neutron Star Merger: GW170817 and GRB~170817A,''
Astrophys.\ J.\ Lett.\ \textbf{848}, L13 (2017).

\bibitem{FermiLAT2009}
Fermi LAT Collaboration, ``A Limit on the Variation of the Speed of
Light Arising from Quantum Gravity Effects,'' Nature \textbf{462},
331 (2009).

\bibitem{CODATA2022}
E.~Tiesinga, P.~J.~Mohr, D.~B.~Newell, and B.~N.~Taylor,
``CODATA Recommended Values of the Fundamental Physical
Constants: 2022,'' Rev.\ Mod.\ Phys.\ \textbf{96}, 025004 (2024).

\bibitem{RS_Axioms}
J.~Washburn, ``The Algebra of Reality: A Recognition Science Derivation
of Physical Law,'' \emph{Axioms} \textbf{15}(2), 90 (2026).
\texttt{doi:10.3390/axioms15020090}

\end{thebibliography}

\end{document}
